\section{Conclusion}
\label{sec:conclusion}

World-model research has so far been organized around individual,
visually rich games, with a large fraction of model capacity spent
reproducing pixels rather than the symbolic mechanics that determine
what happens next. The complementary regime is also worth
investigating: a single dynamics model spanning a family of
mechanically distinct, fully discrete environments, with a substrate
light enough to study at the granularity of individual rules on
commodity hardware. This is largely orthogonal to the dense-pixel
video-world-model line of work, which is expensive to iterate on, and
opens a different research surface, organized around program
structure and rule-system generalization rather than visual fidelity.

Toward this end we proposed a rule-conditioned Neural Cellular
Automaton world model. Three design choices work in concert.
\textbf{(i) An iterated local-update body.} The dynamics body is a
single shared local update rule applied for several iterations on a
hidden grid, mirroring the recursive outer loop of a PuzzleScript
engine that fires a sequence of atomic local rewrites until
convergence. This inductive bias yields parameter-efficient world
models that outperform non-iterative CNN, U-Net, and ViT baselines
on autoregressive rollout fidelity, both for a single rule family
and across many jointly trained games
(\S\ref{sec:results-baselines}, \S\ref{sec:results-scaling}).
\textbf{(ii) Rule conditioning.} Cross-attention to slot embeddings
of each game's tokenized rule text lets a single set of weights
jointly model many games at once, with in-distribution per-cell
fidelity remaining high as the corpus widens
(Table~\ref{tab:indist-scaling}).
\textbf{(iii) Training-data scaling.} Broadening the training
corpus markedly improves out-of-distribution performance on an
unseen-game held-out set, closing most of the gap to the identity
baseline and yielding held-out wins on a non-trivial fraction of
games (\S\ref{sec:results-encoder-ablation}). The headline
rule-conditional contribution to out-of-distribution generalization
is driven primarily by broad-corpus rule diversity rather than the
encoder architecture in isolation, in line with the broader pattern
that generalization is a function of training-distribution coverage.

A general rule-conditioned world model is desirable for several
downstream applications: sample-efficient policy learning and
gradient-based planning across many environments without per-game
retraining, fast adaptation to novel rule systems on interactive
world-model benchmarks where rule inference from limited
interaction is the central
challenge~\cite{warrier2025benchmarking,arcprize2026arcagi3}, and a
bridge between neural world models and symbolic program learning
through a learned latent space that doubles as a search space over
executable game programs. We report preliminary probes of the
latter two in the appendix: random samples from a Gaussian fit to
the empirical slot prior decode to PuzzleScript programs the
original engine compiles and runs
(Appendix~\ref{app:autoencoder}), and freezing the dynamics model
while optimizing the slot from a small tape of an unseen game's
transitions drives per-cell BCE down by orders of magnitude
(Appendix~\ref{app:inverse-fit}). Scaling these probes to the full
held-out corpus and pairing inverse-fit slots with engine-loadable
program reconstructions remain open.

The combination of an efficient substrate, a domain with non-trivial
mechanical complexity, and a learned latent space of embedded games
makes rule-conditioned NCAs a useful sparse, dynamics-first
counterpoint to the dense visual world models that currently dominate
the field, and a natural bridge between neural world models and
symbolic program learning.
